% A.59: registered after full root and DE proof/source/certificate audits PASS.
\subsection{Normalizers for every cover with nilpotent Galois closure}

\begin{statement}{A.59}{Proved}{All-prime normalizer recognition for nilpotent geometric covers}
Let \(U=E\setminus\{O\}\), with \(E:y^2=x^3-x\), and let
\(Y\to U\) be a geometrically connected finite etale cover over
a number field. Suppose that its geometric Galois closure has
finite nilpotent group. Then, for every coefficient prime
\(\ell\), every nontrivial closed procyclic subgroup of
\(\pi_1(Y_{\overline K})^{(\ell)}\) whose arithmetic normalizer
has open image is contained in a conjugate of cusp inertia.
In particular this holds for every cover whose geometric Galois
closure is a group of prime-power order, including non-Galois covers.

More generally, for a fixed coefficient prime \(\ell\), it is
enough to have an actual factorization
\[
                           Y\longrightarrow V\longrightarrow U
                                                        \tag{A.59.1}
\]
in which the geometric Galois closure of the first map is an
\(\ell\)-group and that of the second has order prime to
\(\ell\). All maps and normalizers here are actual arithmetic
fundamental-group maps after finite descent. No full fixed-group
assertion or cofinality among arbitrary covers is claimed.
\end{statement}

\begin{proof}
We first prove the assertion for (A.59.1), then construct this
factorization for every coefficient prime in the nilpotent case.
Finite extensions of the number field are harmless: intersecting
an arithmetic normalizer with the inverse image of an open
Galois subgroup preserves openness of its arithmetic image,
while the geometric inertia subgroups do not change.
We may therefore arrange that the covers, their Galois closures
and all intermediate covers used below are defined over one
finite field extension.

\paragraph{The prime-to-coefficient intermediate curve has good reduction.}
The elliptic curve \(E\) has potentially good reduction at
every finite place, including the places above two: its
\(j\)-invariant is the integer \(1728\), and one may use
\cite[Proposition VII.5.5]{fghSilvermanAEC}. Once its good
model is chosen, the origin is a smooth section, so
\((E,O)\) is a good marked pair.

At a place of residue characteristic \(\ell\), apply the
prime-to-\(\ell\) specialization theorem to this good pair
and the geometric Galois closure of \(V\to U\). After a
finite extension the actual cover extends tamely over the pair;
see \cite[Expos\'e XIII, Corollary 2.9, Section 2.10 and
the proof of Corollary 2.12]{fghSGA1}, with the explicit
comparison already proved in A.56. Its proper source is smooth:
away from the origin the map is etale, and above the origin
its completed local equation after an etale base change is
\(R[[t]]\to R[[z]],\ t=z^e\), with \(e\) prime to
\(\ell\). The same statement holds for subgroup quotients.
Consequently the proper compactification \(\overline V\)
has potentially good \emph{curve} reduction at this place.

A nonconstant map \(\overline V\to E\) forces its genus
to be at least one. If that genus is at least two, A.57
proves its pro-\(\ell\) normalizer criterion. If it is
one, A.40 gives the criterion for the actual isogeny
cover with its removed origin fibre. Thus the normalizer
criterion holds for \(V\) in every case, including
\(\ell=2\). This step needs no restriction on the
degree of \(V\); only its prime-to-\(\ell\) Galois
closure is used.

\paragraph{The first map gives an actual open pro-ell embedding.}
Put \(\Pi=\pi_1(V_{\overline K})\), and let
\(H\subset\Pi\) be the open subgroup corresponding to
\(Y\). Its normal core \(C\) has \(\Pi/C\) an
\(\ell\)-group. The pro-\(\ell\) topology of \(H\)
is induced from \(\Pi\). For completeness, if
\(N\triangleleft H\) is open with \(H/N\) an
\(\ell\)-group, then \(N\cap C\) is normal in
\(C\). Intersect its finitely many conjugates under
\(\Pi/C\), obtaining \(M\triangleleft\Pi\) with
\(M\subset N\). The quotient \(C/M\) embeds in the
product of the \(\ell\)-groups
\(C/g(N\cap C)g^{-1}\). Hence \(\Pi/M\) is an
\(\ell\)-group. This proves the required induced topology
and the continuous open injection
\[
               H^{(\ell)}\lhook\joinrel\longrightarrow
                                  \Pi^{(\ell)}.       \tag{A.59.2}
\]
This is the general completion argument in A.22; it uses
an \(\ell\)-group Galois closure, not merely an
\(\ell\)-power degree.

The actual arithmetic map induced by \(Y\to V\)
has (A.59.2) on its geometric kernels and the identity
on the common field's absolute Galois quotient. It is
therefore itself injective, with open image. In particular
the arithmetic normalizer of a subgroup of \(H^{(\ell)}\)
maps into its arithmetic normalizer in the geometrically
pro-\(\ell\) group of \(V\), with the same open
arithmetic image.

Let \(A\subset H^{(\ell)}\) be a nontrivial closed
procyclic subgroup with that property. By the criterion
for \(V\), its image is contained in some geometric
inertia conjugate \(J\). The intersection
\[
                         H^{(\ell)}\cap J             \tag{A.59.3}
\]
is exactly the full pro-\(\ell\) inertia of a cusp
of the actual curve \(Y\). Indeed the usual geometric
covering correspondence describes the branches above a
cusp by double cosets, and the inertia of the branch
is the corresponding subgroup intersection. Equivalently,
the local map \(t=z^e\) identifies it with
\(e\Z_\ell(1)\subset\Z_\ell(1)\).
Here \(e\) is an \(\ell\)-power, since the Galois
closure of \(Y\to V\) is an \(\ell\)-group. The
induced-topology comparison makes this exactly the
inertia in (A.59.2), with no additional quotient kernel.
Thus \(A\) is cuspidal on \(Y\), proving the general
factorization assertion.

\paragraph{Every nilpotent Galois closure has the required factorization.}
Write \(\pi=\pi_1(U_{\overline K})\), let \(H\subset\pi\)
correspond to the original cover, let \(N\) be its normal
core, and put \(G=\pi/N\), \(A=H/N\). For the chosen
coefficient prime write the finite nilpotent group as
\[
                   G=S\times T,
 \qquad S\text{ an }\ell\text{-group},\quad
              |T|\text{ prime to }\ell.               \tag{A.59.4}
\]
Every subgroup splits accordingly:
\[
             A=A_S\times A_T,
             \quad A_S=A\cap S,\quad A_T=A\cap T.
                                                        \tag{A.59.5}
\]
To see this without a normality assumption on \(A\),
take \((s,t)\in A\) and choose an integer congruent
to one modulo \(|S|\) and zero modulo \(|T|\).
That power of \((s,t)\) is \((s,1)\in A\); the
complementary power gives \((1,t)\in A\). This
proves (A.59.5) for arbitrary subgroups.

Let \(B=S\times A_T\), and let \(V\to U\) be
the actual intermediate cover corresponding to the
inverse image of \(B\). Its geometric Galois closure
has group
\[
 \begin{aligned}
 G/\operatorname{core}_G(B)
    &\simeq T/\operatorname{core}_T(A_T),\\
 B/\operatorname{core}_B(A)
    &\simeq S/\operatorname{core}_S(A_S).
 \end{aligned}                                          \tag{A.59.6}
\]
The first group has order prime to \(\ell\); the
second, which is the Galois closure group of
\(Y\to V\), is an \(\ell\)-group.
The formulas follow directly by taking conjugates
in the direct products; in the second formula
\(A_T\) is normal in the \(T\)-factor \(A_T\)
of \(B\). Thus (A.59.1) has exactly the required
two Galois-closure hypotheses, and the preceding
argument applies.

This proves the all-prime assertion simultaneously
for every cover with finite nilpotent geometric
Galois closure. In particular it gives an independent
normalizer proof for all prime-power Heisenberg
covers previously considered. Their full fixed
subgroups still require separate arguments such
as A.22 and A.41--A.42. The finite nilpotent
quotients are not asserted to be cofinal among
all finite geometric quotients.
The exact certificate \texttt{code/a\_nilpotent\_cover\_normalizers.py}
checks every subgroup and both core formulas in three mixed-prime
examples, including nonnormal subgroups. The general proof above
is independent of that finite sample.
\end{proof}
